\documentclass[11pt,a4paper]{article}

% ------------------------------------------------------------------
% Carlos Iván Pineda Santiago — Currículum (Español)
% Compila con Tectonic / XeLaTeX
% ------------------------------------------------------------------
\usepackage[top=1.5cm,bottom=1.5cm,left=1.7cm,right=1.7cm]{geometry}
\usepackage{xcolor}
\usepackage{titlesec}
\usepackage{enumitem}
\usepackage{tabularx}
\usepackage[hidelinks]{hyperref}

% ---- Colores (blanco y negro) ----
\definecolor{accent}{HTML}{000000}
\definecolor{ink}{HTML}{000000}
\definecolor{muted}{HTML}{555555}

\color{ink}
\hypersetup{colorlinks=true, urlcolor=ink, linkcolor=ink}

% ---- Estilo de secciones ----
\titleformat{\section}
  {\large\bfseries\color{accent}}{}{0em}{}[\vspace{-0.5em}{\color{accent}\rule{\linewidth}{0.9pt}}]
\titlespacing{\section}{0pt}{1.15em}{0.7em}

\setlist[itemize]{leftmargin=1.2em, itemsep=3pt, topsep=3pt, parsep=0pt,
  label={\color{accent}\small\textbullet}}

\setlength{\parindent}{0pt}
\setlength{\parskip}{0pt}
\pagestyle{empty}

% ---- Entrada de dos líneas ----
\newcommand{\cventry}[4]{%
  \begin{tabularx}{\linewidth}{@{}X r@{}}
    \textbf{#1} & {\color{muted}#2}\\[2pt]
    {\color{muted}\itshape #3} & {\color{muted}\itshape #4}\\
  \end{tabularx}%
  \vspace{-0.05em}%
}
\newcommand{\sep}{\hspace{0.4em}\textcolor{muted}{\textbullet}\hspace{0.4em}}

\begin{document}

% ============================ ENCABEZADO ============================
\begin{center}
  {\Huge\bfseries\color{ink} Carlos Iván Pineda Santiago}\\[4pt]
  {\large\color{accent} Ingeniero de Software y Científico de Datos \quad\textbar\quad Fundador de QuantCoder}\\[8pt]
  \small
  +52 55 4254 5603 \sep
  \href{mailto:ivan31416neda@gmail.com}{ivan31416neda@gmail.com} \sep
  \href{https://ivan31416neda.com}{ivan31416neda.com}\\[3pt]
  \href{https://www.linkedin.com/in/carlos-iván-pineda-santiago-325346107/}{linkedin.com/in/carlos-iván-pineda-santiago} \sep
  \href{https://github.com/Ivan252512}{github.com/Ivan252512}
\end{center}

% ============================ PERFIL ============================
\section{Perfil Profesional}
Ingeniero de Software y Científico de Datos con \textbf{más de 5 años} construyendo servicios backend escalables en Python y NodeJS sobre \textbf{AWS}, especializado en \textbf{Fintech} y \textbf{bienes raíces}. Fundador de \textbf{QuantCoder}, una plataforma de trading algorítmico de criptomonedas basada en machine learning. Certificado en Ciencia de Datos y Machine Learning por \textbf{MIT} y la \textbf{UNAM}. Uno la ingeniería de software robusta con el machine learning aplicado para construir y desplegar productos de alto impacto basados en datos, con fuerte interés en el trading algorítmico y la investigación cuantitativa.

% ============================ EXPERIENCIA ============================
\section{Experiencia Profesional}

\cventry{QuantCoder \, — \, Fundador y Propietario}{2023 -- Presente}{Plataforma de trading algorítmico cuantitativo \, \textperiodcentered\ \, \href{https://quantcoderpro.com}{quantcoderpro.com}}{Remoto}
\begin{itemize}
  \item Fundé y construí de extremo a extremo una plataforma que analiza más de 15 pares de criptomonedas (BTC, ETH, SOL, XRP, ADA, DOT\ldots) con modelos de machine learning entrenados con datos de Binance en múltiples temporalidades (4h, 1d, 1 sem, 1 mes), generando señales en tiempo real \textbf{LONG / SHORT / SIDEWAYS}.
  \item Publica automáticamente gráficas de análisis técnico en Facebook y X ante cambios de tendencia, además de alertas personalizadas por Telegram a los suscriptores.
  \item Dashboard web con señales en vivo, gráficas interactivas por par, análisis con GPT-4, un feed diario de noticias financieras en español y suscripciones vía Stripe.
  \item 100\% serverless en AWS (Lambda, DynamoDB, Step Functions, S3, CloudFront) con arquitectura limpia y clasificadores propios Random Forest y K-Means.
\end{itemize}

\cventry{Konfío \, — \, Desarrollador de Software (Finanzas)}{2024 -- Presente}{Ingeniería backend para productos bancarios}{Ciudad de México}
\begin{itemize}
  \item Desarrollo de servicios backend en Python y NodeJS sobre AWS para productos financieros intensivos en datos.
\end{itemize}

\cventry{GBM \, — \, Desarrollador de Software (Finanzas)}{2022 -- 2024}{Plataformas de trading e inversión}{Ciudad de México}
\begin{itemize}
  \item Desarrollé servicios en Python, NodeJS y C\# para plataformas de trading e inversión.
  \item Integré datos de mercado y series de tiempo en APIs usadas por herramientas internas y dashboards.
\end{itemize}

\cventry{TuHabi \, — \, Desarrollador de Software (Bienes Raíces)}{2021 -- 2022}{Servicios backend PropTech}{Ciudad de México}
\begin{itemize}
  \item Implementé servicios backend en Python (Django) y NodeJS sobre AWS usando grandes conjuntos de datos inmobiliarios.
\end{itemize}

\cventry{Codster \, — \, Desarrollador de Software (Consultoría)}{2019 -- 2021}{Software a la medida para diversos clientes}{Ciudad de México}
\begin{itemize}
  \item Entregué soluciones backend a la medida en Python, Django, Flask y NodeJS para diversos clientes.
\end{itemize}

% ============================ PROYECTOS ============================
\section{Proyectos Destacados}

\cventry{Binance Swing Trading Bot}{2021}{Origen de QuantCoder \, \textperiodcentered\ \, Python, ML, AWS}{}
\begin{itemize}
  \item Bot de trading algorítmico para Binance en Python (NumPy, Scikit-learn, PyTorch, TensorFlow), desplegado en AWS.
  \item Experimenté con K-Means, redes neuronales, Random Forest y XGBoost para detección de patrones y generación de señales.
\end{itemize}

\cventry{WSEN Latina \, — \, Sitio Web de Red Científica}{2018}{\href{https://wsen-latina.unam.mx}{wsen-latina.unam.mx} \, \textperiodcentered\ \, UNAM}{}
\begin{itemize}
  \item Sitio web de red científica construido con PHP, Apache y JavaScript.
\end{itemize}

% ============================ HABILIDADES ============================
\section{Habilidades Técnicas}
\begin{itemize}[label={}, leftmargin=0em, itemsep=4pt]
  \item \textbf{\color{accent}Programación:} Python, JavaScript, C\#, SQL
  \item \textbf{\color{accent}Ciencia de Datos / ML:} NumPy, Pandas, Scikit-learn, PyTorch, TensorFlow, XGBoost, Matplotlib / Seaborn
  \item \textbf{\color{accent}Backend y Nube:} Django, Flask, Serverless, React, AWS (Lambda, API Gateway, S3, RDS, DynamoDB), PostgreSQL, MySQL, MongoDB, Docker
  \item \textbf{\color{accent}Temas:} Limpieza de datos, análisis exploratorio, probabilidad y estadística, aprendizaje supervisado y no supervisado, optimización en ML, comunicación de resultados
\end{itemize}

% ============================ FORMACIÓN ============================
\section{Formación Académica}
\cventry{Ingeniería en Sistemas Computacionales}{2019}{Instituto de Estudios Universitarios (IEU)}{}
\vspace{0.5em}
\cventry{Licenciatura en Ciencias de la Tierra (Ciencias Espaciales)}{2018}{Universidad Nacional Autónoma de México (UNAM)}{}

% ============================ CERTIFICACIONES ============================
\section{Certificaciones en Ciencia de Datos}
\cventry{\href{https://olympus1.mygreatlearning.com/certificate/THUSYCVQ}{Ciencia de Datos y Machine Learning: Toma de Decisiones Basada en Datos}}{2024}{MIT Schwarzman College of Computing \, \textperiodcentered\ \, \href{https://olympus1.mygreatlearning.com/certificate/THUSYCVQ}{\underline{verificar}}}{}
\vspace{0.5em}
\cventry{\href{https://ivan31416neda.com/assets/certs/UNAM_Diplomado_Ciencia_de_Datos.pdf}{Introducción Analítica a la Ciencia de Datos (Diploma)}}{2025}{UNAM — Facultad de Ciencias, Educación Continua \, \textperiodcentered\ \, \href{https://ivan31416neda.com/assets/certs/UNAM_Diplomado_Ciencia_de_Datos.pdf}{\underline{ver diploma}}}{}

\end{document}
